
\section{演化式代码与算法发现}
\label{sec:evolutionary-code-algorithm-discovery}

近年来关于AI自进化的文献中出现了一个引人注目的趋同现象：被改进的对象不再仅仅是模型检查点、提示词或策略，而是可执行的制品——它们可以被检查、修改、评估、归档和选择。在本章中，\emph{演化式代码与算法发现}指代这一类方法。该术语涵盖了提示词优化系统（如OPRO）、程序搜索系统（如FunSearch）、代码库级演化智能体（如AlphaEvolve和OpenEvolve），以及智能体级系统（如ADAS、Darwin G{\"o}del Machine、G{\"o}del Agent、Agent Symbolic Learning、Absolute Zero和SelfEvolve）。这些系统在规模和基底上有所不同，但都实例化了同一个闭环：一组候选制品以语言或代码形式表示；生成器提出变体；评估器分配分数；记忆或归档器存储结果；选择规则决定哪些信息应作为下一次提议的条件。

该闭环对于自进化至关重要，因为它为大语言模型（LLM）提供了一种在不改变基础模型权重的情况下改进系统的结构化方式。标准LLM推理是片段式的：模型接收上下文，输出答案，然后忘记这次尝试，除非用户或外部记忆保存了它。演化式代码发现使这些片段得以持久化。一个失败的程序成为反面示例，一个成功的启发式方法成为父代，一个错误追踪成为变异引导，一个基准测试分数成为标量奖励。经过多次迭代，周围的系统可以积累知识，即使LLM本身保持冻结。这就是为什么代码演化在提示工程、AutoML、具有可验证奖励的强化学习和开放式演化计算之间占据桥梁位置。

本章围绕六个问题展开。第一，OPRO如何展示LLM可以通过自然语言历史而非梯度来充当优化器？第二，如何将LLM视为在程序上执行变异、交叉和选择感知重组的演化算子？第三，为什么FunSearch使程序演化成为数学发现的可信路径，其成就应如何与后来AlphaEvolve在矩阵乘法上的结果区分开来？第四，开源系统揭示了哪些关于工程需求、许可和可复现性的信息？第五，智能体自进化系统如何将同一闭环从单个函数扩展到整个智能体架构？第六，目前有哪些基准测试证据，其方法论局限性是什么？

一个有用的抽象是将\emph{内部求解器}与\emph{外部演化过程}分离。内部求解器可以是提示词、Python函数、调度启发式方法、推理轨迹或多智能体工作流。外部过程负责候选生成、评估、归档维护和搜索策略。在经典演化计算中，变异和交叉通常是语法或数值算子。在LLM驱动的演化中，这些算子变成了语义操作：模型可以读取父代程序，推断其意图，识别瓶颈，引入新的算法思想，并编写一个子代——该子代在词元空间中不是一次小编辑，但在设计空间中可能仍然是一个有意义的后代。这种语义变异是关键优势，也是关键风险。它允许向强大思想的不连续跳跃，但也使安全性、信用分配和可复现性变得更加困难。

因此，最强大的系统将神经创造力与硬验证相结合。OPRO在保留任务上评估候选指令。FunSearch在沙箱中执行程序，仅接受通过语法和语义过滤器的候选。AlphaEvolve及其开源后继者依赖自动化评估器获取客观分数。SE-Agent在SWE-bench Verified上评估演化的推理轨迹。Absolute Zero使用代码执行器验证自生成的任务和解决方案。共同的教训是，自进化不仅仅是"要求模型改进自身"。它是评估生态系统的设计，在其中生成的制品可以安全地竞争、失败，并留下塑造未来世代的痕迹。

\subsection{LLM作为优化器：OPRO}
\label{subsec:opro}

OPRO，即通过提示进行优化（Optimization by PROmpting），引入了一种清晰的LLM作为优化器范式的形式化表述\cite{opro2023}。其核心思想看似简单。系统不是手工设计梯度估计器、贝叶斯代理或遗传算子，而是用自然语言表达优化问题。它向优化器LLM提供候选解决方案及其分数的历史记录，要求生成应该表现更好的新候选，用评分器评估这些候选，然后重复该过程。当候选是下游LLM的指令时，OPRO就变成了提示词优化。当候选是数值变量向量或组合问题的路线时，同一外循环就变成了通用的黑盒优化器。

设$\mathcal{X}$为候选空间，$f:\mathcal{X}\rightarrow\mathbb{R}$为黑盒目标函数。在第$t$次迭代中，OPRO维护一个历史记录
\begin{equation}
  \mathcal{H}_t = \{(x_i, f(x_i))\}_{i=1}^{n_t},
\end{equation}
其中$x_i$可以是自然语言指令、参数元组或序列化为文本的组合对象。元提示构造器$M$将历史记录、任务描述$d$、约束$c$和可选示例$e$转换为提示：
\begin{equation}
  m_t = M(\mathcal{H}_t, d, c, e).
\end{equation}
参数为$\theta$的优化器LLM采样一个或多个新候选，
\begin{equation}
  x_{t+1}^{(j)} \sim p_\theta(x \mid m_t, \tau_{\mathrm{opt}}), \qquad j=1,\ldots,k,
\end{equation}
其中$\tau_{\mathrm{opt}}$通常足够高以保持多样性。过滤重复或无效输出后，对候选进行评分，
\begin{equation}
  s_{t+1}^{(j)} = f\bigl(x_{t+1}^{(j)}\bigr),
\end{equation}
历史记录通过添加新对并可选地截断为最具信息量的最高分示例来更新：
\begin{equation}
  \mathcal{H}_{t+1} = \operatorname{Truncate}\left(\mathcal{H}_{t}\cup\{(x_{t+1}^{(j)}, s_{t+1}^{(j)})\}_{j=1}^{k}\right).
\end{equation}
最终输出通常是$x^*=\arg\max_{(x,s)\in\mathcal{H}_T}s$（最大化），或$\arg\min$（损失最小化）。

这种形式化表述之所以重要，是因为优化器不是在$f$的梯度上训练的。它获得的是经验证据的文本追踪，并被要求推断出更好的动作。因此，提示词成为优化器的接口。如果历史记录是有序的、分桶的或带注释的，模型可以利用序数关系。如果提示词包含高频错误，模型可以针对薄弱区域。如果提示词包含约束，模型可以内化可行性。优化算法部分编码在周围的程序中，部分编码在模型关于解释、启发式方法和问题解决模式的预训练分布中。

在提示词优化中，OPRO使用两层架构。外层是提供给优化器LLM的元提示。它包含先前的指令及其分数、声明更高分数更好的陈述，以及以可解析格式生成新指令的请求。内层是提供给评分器LLM的任务提示。对于每个候选指令，评分器LLM解决一组训练示例，系统计算准确率或任务特定分数。本地材料报告了四种指令放置模式：在问题之前、在问题开头、在问题末尾和在答案开头。这很重要，因为提示词优化不仅是对词语的搜索，也是对这些词语与模型生成策略交互位置的搜索。

实现细节展示了自进化的一个通用工程模式。优化器温度较高，鼓励多样化的候选生成，而评分器温度较低，使评估更具确定性。OPRO的默认指令优化循环可以运行约200个搜索步骤，每步生成多个候选指令。元提示可以保留有限数量的顶级指令而非完整历史，既为了适应上下文长度，也为了强调成功模式。分数可以离散化为桶，使优化器不会过拟合可能是噪声的微小分数差异。候选提取使用明确的分隔符，例如对一种模型系列使用\texttt{<INS>}和\texttt{</INS>}，对另一种使用括号约定；无效、重复、过长或明显格式错误的指令在评分前被过滤。

OPRO的数学意义在于它将优化转化为基于评估历史的条件生成。搜索策略不是手工编码的演化规则，但可以模拟演化行为。高分指令充当精英。低分指令充当反面示例。新采样的指令可能改变措辞、重组来自多个父代的思想，或引入从预训练中学到的新启发式方法。选择在历史截断和排名中是显式的。因此，OPRO是从提示工程到演化搜索的极简桥梁：它缺乏种群数据库、岛屿模型或正式的质量-多样性归档，但已经包含了生成--评估--记忆的核心要素。

对于推理任务，OPRO的优化器可以重新发现熟悉的提示模式，如逐步推理、仔细检查或分解。关键不在于每条发现的指令在语义上都是令人惊讶的。关键在于系统能够针对特定任务分布自动发现、排序和重用此类指令。这种任务适应性在AI自进化中尤为重要，因为部署的智能体通常面临的分布与模型的原始训练混合体不同。类似OPRO的外循环可以在不需要梯度更新的情况下调整智能体的操作提示以适应环境。

OPRO还展示了相同的语言接口可以优化非语言变量。在线性回归示例中，候选可以表示为文本对，如$w=15, b=14$，目标是平方误差，
\begin{equation}
  f(w,b) = \lVert y - (Xw+b)\rVert_2^2.
\end{equation}
在旅行商示例中，候选是序列化为轨迹的路线，按旅行长度评分。优化器LLM不接收关于$w$、$b$或路线的导数。它接收的是先前值和目标分数的记录。这暗示了一个广泛的设计原则：当一个领域可以序列化为LLM可解释的表示，并且候选质量可以自动评估时，LLM就可以充当关于改进的黑盒提议分布。

然而，OPRO也暴露了在后续系统中持续存在的局限性。第一，优化器可能被噪声评估误导。如果评分器LLM是随机的或训练集很小，历史可能会奖励虚假指令。第二，历史表示是一个瓶颈。模型只能以适应上下文的内容为条件，因此摘要和选择策略成为优化器的一部分。第三，OPRO的自然语言变异没有单调改进的形式化保证。即使提示词说"生成更好的指令"，采样结果可能更差、冗余或无效。第四，提示词优化可能对用于评分的训练示例过拟合。因此验证间隔和保留集是必不可少的。

对于自进化综述而言，OPRO的主要贡献是概念性的。它表明，如果搜索状态用语言表达且有可靠的分数可用，语言模型就可以被封装为通用优化器。后来的系统用函数、智能体、轨迹和代码库替换了指令，但保留了相同的核心方程：新制品是从以旧制品及其测量后果的归档为条件的模型中采样的。因此，OPRO是从"LLM写答案"到"LLM优化写答案的过程"这一进程的第一步。

\subsection{LLM作为演化算子：变异、交叉与选择}
\label{subsec:llm-evolutionary-operator}

一旦LLM被用于从评分制品历史中生成候选，将其解释为演化算子就是自然的。经典演化算法通常将变异定义为随机扰动，将交叉定义为父代基因组的重组，将选择定义为偏向高适应度个体繁殖的规则。LLM驱动的演化改变了表示和算子。基因组可以是程序、提示词、工作流或推理轨迹。变异可以是阅读错误追踪后提出的语义编辑。交叉可以是代码注释中描述的两种算法思想的综合。选择可以通过归档来实现，向下一次LLM调用暴露成功的父代，同时隐藏或降权失败者。

一般LLM演化循环可以写为如下。设$A_t$为具有适应度$F(a)$和可选行为描述符$b(a)$的制品归档。父代选择选择一个或多个父代，
\begin{equation}
  (a_{p_1},\ldots,a_{p_m})\sim q_t(a\mid A_t),
\end{equation}
其中$q_t$可以依赖于适应度、新颖性、年龄或生态位覆盖。提示构造器将父代代码、评估追踪、失败和指令序列化为$z_t$。LLM采样一个子代，
\begin{equation}
  a_{c}\sim p_\theta(a\mid z_t, \tau, \mathrm{mode}),
\end{equation}
其中\textit{mode}可以请求变异、重组、修复、简化或定向优化。评估器执行或评判$a_c$，产生$F(a_c)$和$b(a_c)$。归档更新规则决定是否插入、丢弃或隔离子代：
\begin{equation}
  A_{t+1}=U(A_t,a_c,F(a_c),b(a_c)).
\end{equation}
这一形式化涵盖了OPRO、FunSearch、AlphaEvolve、OpenEvolve、CodeEvolve和智能体级变体。区别在于$q_t$、提示构造器、评估器和$U$。

LLM变异比词元级变异更强大，因为它可以在改变实现的同时保持意图。例如，模型可以将二次循环替换为堆，引入记忆化，重写启发式优先级函数，或添加防止已知失败模式的保护。这些更改在编辑距离上可能很大，但在概念距离上可能很小。反之，一个微小的文本编辑在语义上可能是灾难性的。这种不匹配意味着LLM系统中的演化距离不应仅通过词元或行来衡量。应通过行为描述符、测试签名、资源概况和谱系关系来衡量。

LLM交叉同样是语义的。经典交叉算子拼接子串或子树。基于模型的算子可以读取两个父代并产生一个连贯的子代，将第一个父代的数据结构与第二个父代的评分规则结合起来。如果各部分不兼容，它也可以拒绝朴素的组合。SE-Agent等系统在轨迹级别使这一点显式化：重组从多条推理路径中提取有用的片段，构建旨在继承其优势的新路径。开源代码演化系统向模型呈现多个"灵感"程序，为模型提供重组的上下文而不强制机械拼接。

选择仍然部分非神经化，且应该保持如此。原因很简单：自进化需要一个在模型对合理文本的偏好之外的接地信号。如果模型同时提议和宣布成功而没有外部验证，循环可能退化为自我确认。因此选择依赖于单元测试、基准测试、数学目标函数、执行追踪或独立控制的评判器。即使使用LLM评判器，鲁棒的系统也偏好低温评估、多样本、校准集或可验证的子组件。评估器是防止演化语言循环变成纯粹修辞的锚点。

AlphaEvolve是LLM作为演化算子范式的最突出大规模示例\cite{alphaevolve2025}。它被描述为一个由Gemini驱动的编码智能体，用于科学和算法发现。其架构结合了一个用于广度的快速模型Gemini Flash和一个用于深度的更强模型Gemini Pro。实际直觉是探索和利用需要不同的模型经济学。广度需要以可接受的成本和延迟获得大量候选想法。深度需要更昂贵的推理来处理困难的改进、抽象或修复。因此，双模型架构创建了跨越模型能力的搜索策略，而不仅仅是跨越程序候选。

AlphaEvolve将FunSearch风格从单个演化函数扩展到更广泛的代码制品。它使用自动化评估器和演化数据库。公开描述强调质量-多样性搜索，特别是MAP-Elites，而非单一的最佳记录链。质量-多样性方法旨在通过在不同行为生态位中保留高质量制品来照亮解空间。在算法发现中，这很重要，因为达到局部最优的第一个解决方案可能不是下一个突破的踏脚石。一个较慢或不够优雅的候选可能包含一个在后续重组后变得有价值的结构思想。

MAP-Elites维护一个按行为描述符索引的归档。设$\mathcal{B}$为离散化描述符空间，$b(a)\in\mathcal{B}$将制品映射到一个单元格。归档在每个单元格中最多存储一个精英：
\begin{equation}
  E[c] = \arg\max_{a: b(a)=c} F(a).
\end{equation}
当子代$a_c$被评估时，如果其单元格为空或改进了该单元格中已有的精英，则插入：
\begin{equation}
  E[b(a_c)] \leftarrow
  \begin{cases}
    a_c, & E[b(a_c)] = \varnothing,\\
    a_c, & F(a_c) > F(E[b(a_c)]),\\
    E[b(a_c)], & \text{其他情况.}
  \end{cases}
\end{equation}
因此归档不是排行榜，而是高性能多样性的地图。

\begin{algorithm}[t]
\caption{用于演化式代码发现的LLM引导MAP-Elites}
\label{alg:llm-map-elites}
\begin{algorithmic}[1]
\REQUIRE 种子程序$P_0$，评估器$\operatorname{Eval}$，描述符函数$b$，LLM集成$\mathcal{M}$，预算$T$
\STATE 用空单元格初始化归档$E$
\FORALL{$p\in P_0$}
  \STATE $(F(p), b(p), r(p)) \leftarrow \operatorname{Eval}(p)$
  \STATE 如果$p$是单元格$b(p)$的精英，则将其插入$E$
\ENDFOR
\FOR{$t=1$ to $T$}
  \STATE 使用适应度、新颖性、年龄和覆盖权重从$E$中选择一个或多个精英$P_t$
  \STATE 构造包含父代代码、分数、描述符、失败和指令的演化提示
  \STATE 选择模型$m_t\in\mathcal{M}$，例如快速模型用于广度或强模型用于深度
  \STATE 采样候选编辑或完整子代程序$C_t\sim m_t(\cdot)$
  \FORALL{$c\in C_t$}
    \STATE 解析、清理和沙箱化$c$
    \STATE $(F(c), b(c), r(c))\leftarrow \operatorname{Eval}(c)$
    \IF{$c$有效且改进了单元格$b(c)$中的精英}
      \STATE $E[b(c)]\leftarrow c$
    \ELSE
      \STATE 存储失败追踪$r(c)$用于诊断或未来提示
    \ENDIF
  \ENDFOR
\ENDFOR
\RETURN 归档$E$和全局最佳精英$\arg\max_{c\in E}F(c)$
\end{algorithmic}
\end{algorithm}

算法框抽象了类似AlphaEvolve的系统，同时也描述了开源实现。几个设计选择隐藏在提示构造器中。提示可以包括按分数排序的父代、简洁的失败分析、禁止回归列表或请求补丁而非完整文件的指令。它可以要求模型保留函数签名、保持在演化区域内或仅添加确定性代码。它还可以使用模型路由：快速模型提出许多粗略编辑，而更强的模型修复、解释或强化有希望的分支。从这个意义上说，提示不仅是输入字符串；它是算子的规格说明。

AlphaEvolve报告的成就说明了为什么该架构引起了关注。本地材料报告它找到了一种48次乘法的$4\times4$复矩阵乘法算法，被描述为56年来对Strassen时代基线的首次改进。它们还报告，在五十多个数学问题中，AlphaEvolve为约75\%的问题重新发现了最先进结果，并为约20\%找到了更好的解决方案。在基础设施应用中，同一来源报告了数据中心调度中恢复全球0.7\%的计算、简化硬件加速器设计以及在特定注意力相关设置中将LLM训练内核加速23\%或32\%等改进。这些数字是领域特定的，应被视为基于评估器优化的示例，而非普遍保证。

双模型架构是一个重要的系统经验。在演化搜索中，模型调用的边际价值取决于其使用位置。早期世代受益于多样性和廉价探索。后期世代可能需要更深入的推理来逃离平台期。系统可以类似于经典演化中的变异率或温度调度来分配模型预算。廉价模型可以用许多变体填充归档；昂贵模型可以在生态位有希望时、失败追踪暗示精确修复时或父代具有高新颖性时被调用。因此自进化变成了关于制品、评估器和模型的资源分配问题。

质量-多样性还重新定义了成功。单一最佳程序对于部署是有用的，但精英归档对于持续自我改进是有用的。如果环境发生变化，非最佳精英可能成为最佳起点。如果后续提示需要不同权衡的示例，归档可以提供它们。如果人类研究人员想要洞见，地图揭示了哪些行为生态位是容易的、困难的、空的或过度代表的。这对于AI自进化尤为相关，因为开放式改进依赖于保留不能立即最大化当前分数的踏脚石。

主要风险是评估器过拟合和规范博弈。如果评估器在狭窄的输入分布上测量速度，演化可能发现脆弱的捷径。如果测试是可见或可推断的，模型可能产生通过测试但未解决预期任务的代码。如果在没有仔细隔离的共享机器上测量性能，噪声可能驱动选择。因此系统需要隐藏测试、随机输入、资源控制、多重指标和审计追踪。它们还需要关于失败候选的策略：失败是有价值的搜索信息，但在提示中包含过多原始失败可能污染上下文并鼓励退化模式。好的归档同时存储成功和摘要化的失败模式。

\subsection{程序演化与数学发现：FunSearch}
\label{subsec:funsearch}

FunSearch是LLM引导的程序演化可以产生数学发现（而非仅仅是更好的基准提示词）的典型示范\cite{funsearch2024}。该系统将预训练LLM与评估器和种群数据库相结合。它不是直接搜索解决方案，而是搜索生成或优先排序解决方案的程序。这一区别至关重要。原始组合对象可能太大、太脆弱或太无结构，无法高效搜索。程序可以编码可重用的启发式方法、构造规则或跨实例规模泛化的优先级函数。LLM对代码和数学的知识变得有用，因为候选制品是可执行的源文本。

FunSearch流水线有三个核心组件：采样器、评估器和程序数据库。数据库生成包含先前程序变体的提示。采样器请求LLM完成被标记为演化的函数的新版本。评估器解析并执行候选，计算分数，并将成功的程序注册回数据库。用户用两个装饰器标记问题：一个函数是运行/评估入口点，另一个函数是要演化的函数。被演化的函数可以是用于贪婪构造数学对象的优先级函数，或求解器内部使用的启发式方法。

本地实现分析报告FunSearch使用具有岛屿、集群和程序的种群结构。岛屿是独立的子种群。集群将跨测试用例具有相同分数签名的程序分组。在集群内，较短的程序更受偏好，应用了偏向简洁启发式的奥卡姆式偏好。定期岛屿重置移除较弱的岛屿并从强奠基者重新启动，这有助于避免过早收敛同时保留成功的发现。这是一个适应LLM生成代码的经典演化思想。

提示设计是FunSearch最重要的贡献之一。提示可以展示一系列版本，如\texttt{priority\_v0}、\texttt{priority\_v1}和一个带有文档字符串说它应该改进先前版本的空\texttt{priority\_v2}。早期版本带有导入和辅助函数，允许LLM看到演化的代码上下文。版本可以排序，使更强的示例出现在后面，利用Transformer上下文中的近期效应。提示不仅仅是请求解决方案；它搭建了一个演化谱系并邀请模型延续它。

FunSearch的评估器刻意严格。它检查生成的代码是否能够解析、在沙箱和超时内运行、不直接调用祖先函数，并返回数值分数。当LLM输出不完整文本时，候选代码可以通过基于AST的恢复进行修剪，但语义准入取决于执行。评估器还记录一个签名：跨测试实例的分数向量。此签名用于聚类，因此具有相同聚合分数但行为不同的两个程序如果其实例级模式不同，就可以被区分。

简化的FunSearch循环可以描述为：
\begin{algorithm}[t]
\caption{FunSearch风格程序演化}
\label{alg:funsearch-style}
\begin{algorithmic}[1]
\REQUIRE 具有\texttt{@run}评估器和\texttt{@evolve}目标的初始程序；LLM；程序数据库$D$；评估器套件$E_{\mathrm{eval}}$
\STATE 用种子程序在$D$中初始化岛屿
\WHILE{预算剩余}
  \STATE 采样一个岛屿并从高分程序版本构造提示
  \STATE 请求LLM完成下一个演化函数版本
  \STATE 解析候选代码；修剪或拒绝格式错误的主体
  \STATE 在沙箱中对基准实例执行候选
  \IF{候选运行、避免祖先调用并返回数值分数}
    \STATE 计算签名向量和归约分数
    \STATE 将程序注册到相应的岛屿和集群
  \ELSE
    \STATE 丢弃或记录失败
  \ENDIF
  \IF{重置周期已到}
    \STATE 从强岛屿中采样的奠基者重新初始化弱岛屿
  \ENDIF
\ENDWHILE
\RETURN 最佳程序及其数学构造
\end{algorithmic}
\end{algorithm}

FunSearch的头条成就包括cap set问题的新构造和bin packing的更好启发式方法。它还在相关组合设置中发现了解决方案，如admissible sets和循环图独立集。这些结果具有科学意义，因为输出程序可以被检查。人类可以阅读演化后的优先级函数，推理其为何有效，并可能将其转化为数学洞见。这种可检查性使程序搜索区别于黑盒神经预测。

重要的是纠正一个常见的归属混淆。$4\times4$复矩阵乘法对Strassen时代结果的首次改进在本地的AlphaEvolve材料中与AlphaEvolve相关联，而非原始FunSearch论文。FunSearch是确立了用于数学发现的LLM加评估器程序演化的先驱。AlphaEvolve后来将该方法扩展到更大的代码制品并报告了矩阵乘法改进。为历史准确性，表~\ref{tab:funsearch-alphaevolve-results}将FunSearch的发现与AlphaEvolve的后继结果分开列出。

\begin{table*}[t]
\centering
\caption{代表性程序演化发现。矩阵乘法行被纳入是因为它常与FunSearch一起讨论，但报告的结果属于AlphaEvolve而非原始FunSearch系统。}
\label{tab:funsearch-alphaevolve-results}
{\TableTight
\begin{tabularx}{\textwidth}{@{}L{0.16\textwidth}L{0.19\textwidth}L{0.20\textwidth}Y@{}}
\toprule
\textbf{系统} & \textbf{领域} & \textbf{报告结果} & \textbf{解读} \\
\midrule
FunSearch & Cap set构造 & 在所研究的设置中优于先前已知的构造 & 证明演化程序可以产生新的数学对象，而不仅仅是解决训练实例。 \\
FunSearch & Bin packing & 改进的启发式规则 & 表明程序演化可以发现具有可解释逻辑的实用算法。 \\
FunSearch & 循环图和相关组合学 & 在选定问题中的新构造程序 & 支持函数空间搜索在具有评估器的数学领域间泛化的论断。 \\
AlphaEvolve & $4\times4$复矩阵乘法 & 48次乘法，被报告为56年来对Strassen时代基线的首次改进 & 后继系统将程序演化扩展到更广泛的算法代码和质量-多样性搜索。 \\
AlphaEvolve & 50+数学问题 & 本地材料中恢复了约75\%的最先进结果并改进了约20\% & 表明在评估器设计可用时的广度，但结果取决于问题和预算。 \\
\bottomrule
\end{tabularx}}
\end{table*}

方法论教训是，当三个条件对齐时数学发现成为可能。第一，领域必须允许紧凑的可执行表示。第二，评估器必须比发现的价值更便宜且足够可靠以引导搜索。第三，搜索系统必须保持多样性，因为数学突破通常需要踏脚石。FunSearch在多个组合问题上满足这些条件。它不能解决评估模糊、昂贵或依赖于人类品味的领域。

FunSearch还揭示了为什么程序演化不同于普通代码生成。代码生成基准通常要求模型一次性解决问题。程序演化要求模型参与迭代种群过程。模型不仅因第一次采样而被评判；它还因采样能否被选择、通过上下文重组并在世代间改进而被评判。平庸的采样如果引入了后来变体可以利用的结构基序，可能仍然有价值。因此科学生产力的单位是搜索过程，而不是单个补全。

对于AI自进化，FunSearch提供了一个可以从函数提升到智能体的模板。用规划策略、工具选择规则或记忆更新函数替换优先级函数。用任务基准替换数学评估器。用智能体行为签名替换岛屿集群。同一循环就变成了智能体演化。难度增加了，因为智能体具有更大的状态空间、更多的不确定性来源和更多的机会利用评估器漏洞。尽管如此，FunSearch架构仍然是一个基础模式：演化可执行制品，自动评估它们，保持种群，并暴露可解释的谱系。

\subsection{开源实现：OpenEvolve、FunSearch、CodeEvolve和SE-Agent}
\label{subsec:open-source-implementations}

开源实现之所以重要，是因为它们揭示了研究论文背后隐藏的工程面。论文可能声明评估器执行程序，但可重用的系统必须决定如何隔离文件系统、限制内存、解析模型输出、检查点种群、在模型提供商之间路由、从失败的子进程恢复以及向用户呈现结果。它还必须选择许可证和可复现性姿态。在自进化中，这些工程细节不是附带性的；它们定义了什么样的制品可以安全地演化。

FunSearch发布的代码最好被视为研究框架。它暴露了采样器、评估器、程序数据库、装饰器、AST工具和岛屿模型，但将LLM接口和安全执行环境留作用户必须实现或适配的组件。这对于旨在展示科学方法的Nature论文来说是合适的。代码使算法可理解和可审计，但它不是一个用于任意代码库演化的交钥匙平台。

OpenEvolve将自身定位为受AlphaEvolve启发的开源实现。当前代码库将其描述为一种演化编码智能体，将LLM转化为自主代码优化器，支持Python库使用，与OpenAI兼容的提供商一起工作，并包含MAP-Elites质量-多样性、基于岛屿的种群、LLM集成和用于错误反馈的制品侧通道。其README报告了诸如圆填充、自适应排序、滤波器设计、提示词演化和GPU相关优化等示例，代码库采用Apache-2.0许可证。这将OpenEvolve定位为研究级程序演化与面向用户的自动化优化之间的实用桥梁。

CodeEvolve是另一个开源演化编码智能体。其arXiv摘要描述了一种基于岛屿的遗传算法，具有模块化LLM编排、执行反馈、任务特定指标、上下文感知重组、自适应元提示和定向优化。代码库README描述了Apache-2.0许可、可选MAP-Elites集成、分布式岛屿、提示种群、SEARCH/REPLACE差异应用、资源受限沙箱评估以及通过OpenAI兼容API支持开放权重和专有模型。因此CodeEvolve作为多岛屿编排、探索-利用调度和结构化补丁生成的工程参考尤其有用。

SE-Agent在基底上有所不同。它演化多步LLM代码智能体的推理轨迹，而非直接仅演化源函数。其代码库描述了三种操作：修订、重组和精炼。修订分析失败的轨迹并生成替代策略。重组从多条路径的优势中合成新轨迹。精炼从有希望的路径中移除冗余和风险模式。代码库报告在SWE-bench Verified上达到80\%的Top-1性能，并采用MIT许可证。这使SE-Agent成为演化思想应用于智能体轨迹和问题解决过程（而非仅代码制品）的关键示例。

\begin{table*}[t]
\centering
\caption{代表性开放和半开放演化代码或轨迹系统比较。公共代码库细节可能变化；本表记录了本地研究材料和为本草案执行的Web验证中可用的证据。}
\label{tab:open-source-evolution-systems}
{\TableTight
\begin{tabularx}{\textwidth}{@{}L{0.12\textwidth}L{0.18\textwidth}L{0.34\textwidth}L{0.14\textwidth}Y@{}}
\toprule
\textbf{系统} & \textbf{演化制品} & \textbf{搜索与评估机制} & \textbf{状态} & \textbf{最佳用例} \\
\midrule
FunSearch & 问题脚手架内标记的单个Python函数 & 岛屿模型、按分数签名的程序集群、LLM补全版本化函数，以及用于cap set、bin packing、admissible/循环图任务的沙箱执行。 & Apache-2.0代码；本地笔记中为CC-BY材料。 & 可解释的数学程序发现。 \\
OpenEvolve & 函数、程序、提示词和优化目标 & MAP-Elites质量-多样性、岛屿、LLM集成、制品反馈，以及用户定义的评估器用于圆填充、提示词优化、排序和科学计算。 & Apache-2.0代码库。 & 具有提供商灵活性的实用AlphaEvolve风格实验。 \\
CodeEvolve & 代码区域和算法解决方案 & 基于岛屿的遗传算法、上下文感知重组、自适应元提示、可选MAP-Elites和沙箱化任务特定指标。 & Apache-2.0代码库。 & 可复现的多岛屿代码演化和消融研究。 \\
SE-Agent & 代码智能体的多步推理轨迹 & 跨轨迹的修订、重组和精炼，在SWE-bench Verified上评估，代码库报告80\% Top-1。 & MIT代码库。 & 改进真实软件问题上的智能体推理过程。 \\
AlphaEvolve & 更广泛的算法代码制品 & Gemini Flash/Pro双模型演化，具有用于数学问题、矩阵乘法、数据中心调度、芯片设计和LLM训练内核的质量-多样性归档。 & 本地笔记中未普遍公开。 & 大规模评估器丰富的算法发现。 \\
\bottomrule
\end{tabularx}}
\end{table*}

比较突出了三个轴。第一个轴是制品粒度。FunSearch演化一个函数；OpenEvolve和CodeEvolve可以演化更大的代码区域；SE-Agent演化轨迹；AlphaEvolve针对更广泛的算法代码。更大的粒度增加了表达能力但也增加了失败模式。单个函数可以相对干净地沙箱化和评分。代码库级智能体可能需要依赖安装、编译、集成测试、性能分析和安全隔离。

第二个轴是多样性管理。FunSearch使用岛屿和签名集群。OpenEvolve和CodeEvolve暴露MAP-Elites或基于岛屿的种群。SE-Agent使用轨迹池和重组。这些机制都解决同一个问题：贪心的仅最佳搜索会丢失踏脚石。在自进化系统中，多样性不是装饰性的。它是关于如何解决任务的替代假设的记忆。

第三个轴是评估。FunSearch和AlphaEvolve在评估器是清晰的目标函数时最强。OpenEvolve和CodeEvolve通过让用户定义评估器来泛化，增加了适用性但将责任转移给从业者。SE-Agent的评估器是一个软件工程基准，其成功取决于解决真实代码库问题。这更现实但也更昂贵且更难复现。制品越开放，评估就越成为系统工程问题。

许可也影响采用。FunSearch、OpenEvolve和CodeEvolve中的Apache-2.0代码支持带专利授予的宽松使用，而SE-Agent的MIT许可同样宽松但更短。相比之下，AlphaEvolve在论文和报告中描述，但在本地笔记中作为完整公开系统并不普遍可用。这意味着开源后继者不仅仅是克隆；它们是社区用来测试AlphaEvolve/FunSearch范式的哪些部分可以在大型工业实验室之外复现的制品。

一个实用要点是开源演化系统需要鲁棒的默认值。用户不应该从头设计如何检查点归档、如何摘要失败追踪、如何分离训练和验证评估或如何防止演化程序读取隐藏答案。下一代框架应该将这些视为一等关注点。配置应该暴露模型路由、评估器预算、归档策略、多样性描述符、沙箱限制和验证门控。没有这些控制，自进化可能在演示中看似有效但在科学复现中失败。

\subsection{智能体自进化系统}
\label{subsec:agent-self-evolution-systems}

上述系统演化提示词、函数、代码或轨迹。智能体自进化进一步提高了抽象级别：演化制品是智能体架构、工作流、工具接口、记忆系统或智能体自身的源代码。这种转变在概念上很重要，因为智能体不是单一的输入-输出函数。它们执行多步推理、调用工具、维护状态、与环境交互，并可能与其他智能体协调。因此演化不仅要优化智能体说什么，还要优化它如何感知、规划、行动、记忆、批评和修改自身。

ADAS，即自动化智能体系统设计，形式化了这一方向\cite{adas2025}。其元智能体搜索算法使用元智能体编写新智能体架构的代码。搜索空间是实现智能体的有效Python程序集合，原则上图灵完备。每次迭代中，元智能体回顾先前发现的设计及其性能，编写新智能体，在基准任务上评估它，并记录结果。输出是归档中发现的最佳设计。这是应用于智能体设计而非数学启发式的演化代码发现。

ADAS的重要性在于其搜索空间的范围。手工设计的智能体框架通常暴露固定菜单：思维链、反思、工具使用、辩论、树搜索或规划。代码级设计搜索可以组合这些基序或发明新的控制流。本地材料报告ADAS发现的智能体在ARC和Google Research Football等基准上优于最先进的手工设计智能体，并跨领域和跨模型骨干迁移。确切的数字在本地摘要中不如其他一些系统详细，但定性结果是智能体设计本身成为优化目标。

Darwin G{\"o}del Machine（DGM）通过开放式归档和自我修改扩展了这一思想\cite{darwingodelmachine2025}。它维护一个不断增长的智能体实现归档，以多样性偏向采样父代智能体，使用基础模型修改智能体代码，在编码基准上评估子代，并保留改进。本地材料报告SWE-bench从20.0\%提升到50.0\%，Polyglot多语言编码基准从14.2\%提升到30.7\%。DGM之所以是"Darwin"，是因为它使用种群搜索和选择；之所以是"G{\"o}del"，是因为智能体可以在受自指改进启发的开放过程中修改自己的代码。

DGM说明了归档相对于单谱系的价值。如果子代失败，父代仍然保留。如果一个分支发现了有用的工具使用模式但不是最佳分数，它可以作为踏脚石持续存在。如果后续任务分布变化，归档可能包含比当前冠军更好的起点。这比爬山更鲁棒，更接近开放式演化。权衡是成本：每个子代必须在有意义的软件任务上评估，代码修改可能破坏语法、依赖或安全假设。

G{\"o}del Agent采取了更直接的自指路线\cite{yin2024godel}。它允许智能体使用LLM驱动的猴子补丁检查和修改自己的运行时逻辑。循环是：在任务上执行，收集结果和反馈，分析当前代码和行为，提出补丁，在运行时应用它，验证，然后接受或回滚。这种架构更接近在线自我修改而非种群级演化。其安全机制包括回滚、验证和高级目标约束。本地材料报告在HotPotQA、ALFWorld和更广泛的文本任务上的改进，尽管摘要的数字是定性的而非详细表格。

G{\"o}del Agent暴露了不同的风险特征。运行时猴子补丁灵活且快速，但它是Python特定的并且可能引入不确定性。它还模糊了被评估候选与评估器环境之间的边界。像DGM这样的种群归档可以隔离子代；运行时补丁智能体必须特别小心回滚和验证，因为正在执行的系统就是被更改的对象。对于安全关键的自进化，这种区别很重要。

Agent Symbolic Learning提供了智能体自进化的更结构化和更少代码破坏性的视角\cite{agentsymboliclearning2024}。它将智能体视为符号网络，其可学习的"权重"是提示词、工具和流水线连接。执行任务后，LLM从轨迹中产生语言损失。然后语言梯度通过符号网络反向传播，优化器更新提示节点、工具节点或流水线节点。本地材料中的公式为：
\begin{equation}
  \mathcal{L}_{\mathrm{lang}} = \operatorname{LLM}(\mathcal{P}_{\mathrm{loss}}(\tau)),
\end{equation}
\begin{equation}
  \nabla^n_{\mathrm{lang}} = \operatorname{LLM}\left(\mathcal{P}_{\mathrm{gradient}}(\nabla^{n+1}_{\mathrm{lang}}, \mathcal{I}_n, \mathcal{O}_n, \mathcal{P}_n, \mathcal{T}_n, \mathcal{L}_{\mathrm{lang}})\right),
\end{equation}
随后对提示词、工具和流水线结构进行符号更新。隐喻是语言空间中的反向传播。与DGM不同，更新目标不是任意源代码；它是结构化的智能体图。

本地材料中的基准数字表明广泛但适度的增益。在HotPotQA上，符号学习将F1从36.2提高到39.1，精确匹配从22.8提高到25.6。在使用GPT-4骨干的MATH上，准确率从56.0\%上升到60.7\%。在HumanEval上，pass@1从68.3\%提高到73.2\%。在1到4的软件开发可执行性评分上，报告的平均值从智能体的2.4增加到3.8。这些结果支持语言空间信用分配可以在不需要权重更新的情况下改进已部署智能体的观点。

Absolute Zero将自进化带入具有可验证奖励且无外部数据的强化学习\cite{absolutezero2025}。单个模型扮演两个角色：任务提议者和任务求解器。它生成推理任务，尝试解决它们，并从代码执行器获得奖励。提议者因任务可解、有挑战性和多样而获得奖励；求解者因正确解决方案而获得奖励。本地材料将此描述为零数据自我对弈RLVR，使用PPO或GRPO风格的更新用于两个角色。关键方程是双奖励循环：
\begin{equation}
  t \sim \pi_{\mathrm{prop}}(\cdot\mid s),\qquad y\sim\pi_{\mathrm{solve}}(\cdot\mid t),\qquad r=\operatorname{Execute}(y,t).
\end{equation}
模型通过生成自己的课程来改进。

Absolute Zero在狭义上不是演化代码搜索，但它属于本章，因为它共享自进化循环：生成任务，评估结果，更新生成器/求解器，然后重复。它还展示了可验证性的重要性。没有代码执行器，自生成的任务可能漂移到模糊或自利的问题。有了基于执行的奖励，系统可以在自动接地的课程上训练。本地材料报告AZR-Coder-7B在7B模型中达到编码平均的最先进结果，并在HumanEval+、MBPP+、LiveCodeBench、GSM8K和MATH上保持竞争力，尽管使用零外部训练数据。

SelfEvolve，来自2023年，是一个早期的代码自我改进流水线\cite{selfevolve2023}。它有两个阶段。首先，LLM从自身参数生成知识，如API文档或算法描述。其次，生成的代码在沙箱中执行，错误消息被反馈用于迭代修订。本地材料给出了知识条件生成和修订的公式，包括：
\begin{equation}
  P(Y\mid K)=\prod_i p_\theta(Y_i\mid Y_{<i},X,K),
\end{equation}
其中$K$是自生成的知识，以及以原始问题、生成代码、知识和错误消息为条件的修订分布。SelfEvolve更接近自调试而非种群演化，但它提供了后来系统重用的关键机制：执行反馈作为下一次尝试的自然语言训练信号。

综合来看，这些系统定义了一个谱系。ADAS在代码空间中搜索智能体架构。DGM演化自修改编码智能体的种群。G{\"o}del Agent执行运行时自我修改。Agent Symbolic Learning通过语言梯度更新结构化提示词、工具和流水线。Absolute Zero通过自生成可验证任务训练模型。SelfEvolve使用自生成知识和执行错误迭代优化代码。所有这些都实例化了自进化，但它们的更新基底不同：架构代码、智能体源码、运行时方法、符号图权重、模型策略和生成代码。

\begin{table*}[t]
\centering
\caption{智能体自进化系统按基底、反馈信号和报告的基准证据。}
\label{tab:agent-self-evolution-systems}
{\TableTight
\begin{tabularx}{\textwidth}{@{}L{0.13\textwidth}L{0.18\textwidth}L{0.19\textwidth}L{0.27\textwidth}Y@{}}
\toprule
\textbf{系统} & \textbf{演化基底} & \textbf{反馈信号} & \textbf{本地材料中报告的证据} & \textbf{主要局限} \\
\midrule
ADAS & 实现智能体架构的Python程序 & 智能体任务的基准分数 & 在ARC和GFootball上优于手工设计智能体；跨领域和模型迁移 & 巨大代码空间上的昂贵搜索 \\
DGM & 编码智能体实现的归档 & SWE-bench和Polyglot分数 & SWE-bench 20.0\% $\rightarrow$ 50.0\%；Polyglot 14.2\% $\rightarrow$ 30.7\% & 评估成本和安全自修改 \\
G{\"o}del Agent & 通过猴子补丁的运行时智能体逻辑 & 任务反馈加验证 & 报告在HotPotQA、ALFWorld和多领域文本任务上的增益 & 运行时补丁风险和较弱的可移植性 \\
Agent Symbolic Learning & 提示词、工具和流水线连接 & 语言损失和语言梯度 & HotPotQA F1 36.2 $\rightarrow$ 39.1；MATH 56.0\% $\rightarrow$ 60.7\%；HumanEval 68.3\% $\rightarrow$ 73.2\% & 依赖强LLM生成梯度质量 \\
Absolute Zero & 通过RLVR的提议者/求解器策略 & 自生成任务上的代码执行器奖励 & 本地笔记中7B级编码平均SOTA；数学迁移竞争力 & 课程崩溃和执行器受限领域 \\
SelfEvolve & 生成代码加自生成知识 & 沙箱错误和测试反馈 & DS-1000 49.4 $\rightarrow$ 57.2；HumanEval pass@1 74.39 $\rightarrow$ 85.98 & 主要处理API/断言错误；手写提示词 \\
SE-Agent & 代码智能体的推理轨迹 & SWE-bench Verified结果 & 本地笔记中高达55\%的相对改进；代码库报告80\% Top-1 & 昂贵的多轨迹搜索 \\
\bottomrule
\end{tabularx}}
\end{table*}

一个统一视角是智能体自进化需要三种记忆。第一种是情景记忆：在特定任务尝试中发生了什么。第二种是设计记忆：哪些提示词、工具、代码更改或轨迹在跨尝试中有效。第三种是评估记忆：哪些测量是可信的以及在什么条件下。缺乏情景记忆的系统会重复错误。缺乏设计记忆的系统无法积累改进。缺乏评估记忆的系统会过拟合、退化或奖励黑客。

另一个统一问题是遏制。提示词和符号更新相对容易沙箱化，因为它们改变文本字段。代码更新需要测试、权限、依赖控制和回滚。运行时自我修改需要更强的遏制，因为更改的主体和机制重叠。RL自我对弈需要奖励遏制，因为模型可能学会提出利用验证器的任务。因此自进化既是算法问题，也是基础设施问题。

\subsection{基准比较与证据质量}
\label{subsec:benchmark-comparison}

该领域的基准比较是困难的，因为系统在不同的预算下优化不同的基底。OPRO优化提示词和黑盒变量。FunSearch优化数学程序。AlphaEvolve优化算法和基础设施代码。DGM和SE-Agent针对软件工程智能体。Agent Symbolic Learning跨问答、数学、代码、软件开发和创意写作更新智能体组件。Absolute Zero通过自我对弈训练模型。单一排行榜会模糊这些差异。更好的比较应问：制品是什么？评估器是什么？使用了什么基线？消耗了多少搜索或训练预算？结果是否保留、可复现且抗评估器博弈？

{\TableTight
\begin{longtable}{@{}L{0.13\textwidth}L{0.18\textwidth}L{0.17\textwidth}L{0.18\textwidth}L{0.28\textwidth}@{}}
\caption{本地研究材料和经Web验证的开放代码库中所选数值结果。值应在行内比较，而非跨不相关基准。}
\label{tab:numeric-benchmark-comparison}\\
\toprule
\textbf{系统} & \textbf{基准或领域} & \textbf{基线} & \textbf{报告的演化结果} & \textbf{备注} \\
\midrule
\endfirsthead
\caption[]{所选数值基准结果（续）。}\\
\toprule
\textbf{系统} & \textbf{基准或领域} & \textbf{基线} & \textbf{报告的演化结果} & \textbf{备注} \\
\midrule
\endhead
\bottomrule
\endlastfoot
OPRO & 推理任务上的指令优化 & 初始或先前提示词 & 通过迭代搜索改进的提示词分数 & 确切数字取决于数据集和提示词放置；关键证据是优化轨迹。 \\
OPRO & 线性回归/TSP优化 & 手工或启发式初始候选 & 通过文本条件提议获得的更低目标值 & 展示了超越提示词的黑盒优化。 \\
FunSearch & Cap set & 所研究设置中的先前构造 & 更好的演化构造 & 通过可执行程序搜索的科学发现。 \\
FunSearch & Bin packing & 现有启发式基线 & 改进的启发式程序 & 展示了实用的启发式发现。 \\
AlphaEvolve & $4\times4$复矩阵乘法 & Strassen时代乘法计数 & 48次乘法 & 本地材料中报告为56年来首次改进。 \\
AlphaEvolve & 50+数学问题 & 已知SOTA & 恢复了75\%的SOTA，改进了20\% & 表明了广泛的评估器驱动搜索。 \\
AlphaEvolve & 数据中心调度 & 现有Borg调度策略 & 恢复全球0.7\%计算 & 工业结果；评估器和部署细节是领域特定的。 \\
AlphaEvolve & LLM训练内核 & 现有内核 & 报告的注意力设置中23\%和32\%加速 & 性能取决于硬件和工作负载。 \\
DGM & SWE-bench & 20.0\% & 50.0\% & 基于归档的自改进编码智能体。 \\
DGM & Polyglot编码 & 14.2\% & 30.7\% & 多语言迁移证据。 \\
Agent Symbolic Learning & HotPotQA & F1 36.2 / EM 22.8 & F1 39.1 / EM 25.6 & 结构化语言梯度更新。 \\
Agent Symbolic Learning & MATH & 56.0\% & 60.7\% & 本地材料中的GPT-4骨干。 \\
Agent Symbolic Learning & HumanEval & 68.3\% pass@1 & 73.2\% pass@1 & 无需权重更新的代码生成改进。 \\
SelfEvolve & DS-1000 & ChatGPT 49.4 & 57.2 & 完整流水线，本地笔记中+15.8\%相对。 \\
SelfEvolve & HumanEval & ChatGPT 74.39 pass@1 & 85.98 pass@1 & 接近本地表格中的GPT-4基线。 \\
SelfEvolve & TransCoder & ChatGPT 88.3准确率 / 88.5 pass@1 & 90.4准确率 / 90.9 pass@1 & 翻译和执行精炼。 \\
Absolute Zero & 7B编码平均 & 数据训练的比较方法 & 本地笔记中7B模型SOTA & 零外部数据；确切聚合取决于基准组合。 \\
SE-Agent & SWE-bench Verified & 基础代码智能体 & 高达55\%相对改进；代码库报告80\% Top-1 & 轨迹级演化；经Web验证的代码库声明。 \\
OpenEvolve & 提示词优化示例 & 基线提示词 & 代码库报告示例中+23\% HotpotQA准确率 & 开源框架结果，非通用基准。 \\
CodeEvolve & 算法发现任务 & 开放/封闭基线 & arXiv摘要报告多个任务上的SOTA & 需要论文特定的表格获取每任务的确切值。 \\
\end{longtable}}

该表格有意混合了精确和定性条目，因为源材料在粒度上有所不同。对于某些系统，本地笔记提供了精确的前后值。对于其他系统，它们提供了如"竞争力"、"SOTA"或"优于手工设计智能体"等声明，而没有逐行的数值细分。严格的元分析应回到原始论文，提取任务定义、预算、种子和置信区间，并分离训练、验证和测试性能。此处的目标是绘制证据全景，而非宣布通用赢家。

几个基准模式反复出现。第一，最大的增益通常出现在基线是固定智能体或固定提示词而演化系统可以利用任务特定反馈时。这并不削弱结果；它澄清了自进化是一种适应形式。第二，代码和数学任务过度代表，因为它们具有可验证的奖励。第三，SWE-bench等软件工程基准很有吸引力，因为它们测试现实的智能体行为，但它们成本高且对线束细节敏感。第四，基础设施优化结果即使百分比改进看起来很小也可能具有高度价值。全球范围内0.7\%的计算恢复可能比玩具任务上的大相对增益在经济上更重要。

证据质量取决于评估器独立性。如果候选生成看到确切测试，演化可能过拟合。如果验证在与选择相同的分布上执行，性能可能在分布外退化。如果评估器是LLM评判器，模型偏差和提示词敏感性进入循环。如果基准是公开且广泛使用的，数据污染是可能的。因此，高质量的自进化证据应报告隐藏测试、消融、计算预算、随机种子、失败运行率和回归检查。

消融尤为重要。在演化系统中，改进可以来自许多来源：更多模型调用、更好的提示词、更强的基础模型、更大的归档、更智能的选择、更多样的父代、更好的评估器或仅仅是更多重试。没有消融，很难知道提议的演化机制是否必要。例如，MAP-Elites归档应在相同评估预算下与仅最佳选择进行比较。双模型架构应在相同成本下与单个强模型进行比较。轨迹重组算子应与独立重试和自我反思进行比较。文献开始包含此类消融，但综述读者应将头条分数视为不完整的。

另一个问题是预算归一化。评估一万个候选的系统可能因与算法设计无关的原因优于评估一百个候选的系统。如果用例允许该预算，这不是缺陷，但必须报告。有用的指标包括分数与评估次数、分数与挂钟时间、分数与词元成本和分数与人工工程时间。在工业设置中，相关指标可能是计算投资回报。在科学发现中，相关指标可能是每次评估器调用的新颖性。

对于智能体自进化，迁移是一个关键的基准维度。如果演化智能体仅在选择使用的任务上改进，它是一个专门的优化器。如果其设计迁移到新任务、模型或领域，它更接近通用的自进化方法。本地ADAS笔记强调跨领域和跨模型迁移。DGM报告跨编码基准的迁移。Absolute Zero报告从代码自我对弈到数学推理的跨领域迁移。这些结果很重要，因为自进化理想情况下应该发现可重用机制，而不仅仅是基准特定的技巧。

\subsection{鲁棒AI自进化的设计模式}
\label{subsec:evolution-design-patterns}

在综述的系统中，出现了几种设计模式。第一种是\textbf{表示规范}。演化制品应有清晰的边界。FunSearch标记一个函数。CodeEvolve用结构化差异在标记之间修改代码。Agent Symbolic Learning更新命名的提示词、工具和流水线节点。SE-Agent操作轨迹。清晰的边界使评估、回滚和归因成为可能。

第二种模式是\textbf{验证器优先设计}。评估器应在大规模生成开始之前设计。弱评估器使系统优化错误目标。慢评估器使搜索不切实际。非确定性评估器使选择有噪声。泄露的评估器招致奖励黑客。最成功的系统是那些可以自动且足够廉价地测量候选质量以支持多次迭代的系统。

第三种模式是\textbf{归档优于记忆}。聊天记录不是归档。归档以结构化方式存储制品、分数、描述符、来源、失败追踪和选择状态。它支持父代采样、多样性分析、回滚和审计。OPRO的历史是最小归档。FunSearch的程序数据库是更丰富的归档。DGM的开放式智能体集合和AlphaEvolve的质量-多样性地图更加强大。

第四种模式是\textbf{多样性保持}。仅最佳循环简单但脆弱。它们丢弃了后来可能变得有用的变体。岛屿模型、MAP-Elites、签名聚类和轨迹池都保持替代假设。多样性不仅关乎避免局部最优；它关乎保持可解释性。多样化的归档让研究人员比较解族并识别与成功相关的机制。

第五种模式是\textbf{失败利用}。在普通流水线中，失败的生成被丢弃。在自进化中，失败可以压缩为有用信号：语法错误、超时原因、失败测试名称、反例、性能概况或推理盲点。SelfEvolve使用错误消息进行修订。OpenEvolve风格的制品侧通道将诊断反馈给LLM。SE-Agent使用失败轨迹进行修订。挑战是在不淹没上下文的情况下摘要失败。

第六种模式是\textbf{模型预算调度}。并非每次生成都值得最强的模型。AlphaEvolve的Flash/Pro分割暗示了一个通用方案：使用更便宜的模型或更高温度进行广泛探索，使用更强的模型进行深度修复或利用，并基于归档状态路由调用。开源系统可以用开放权重和专有模型实现类似策略。预算感知的演化对于超越演示的规模化至关重要。

第七种模式是\textbf{回归控制}。任何自进化系统都可能遗忘或退化。代码更改可能改进一个指标同时破坏另一个。提示词更改可能过拟合。智能体补丁可能改变安全行为。鲁棒系统需要基线、验证门控、金丝雀任务、隐藏测试和回滚。当部署需要可靠性而非仅仅最大基准分数时，接受标准应该是多目标的。

第八种模式是\textbf{人类可读谱系}。当人类可以检查解决方案如何演化时，科学和工程价值增加。FunSearch的版本化函数、DGM的归档和CodeEvolve的谱系追踪都支持谱系分析。谱系可以揭示结果是通过渐进优化、突然引入的想法还是评估器利用产生的。这对于信任和将演化制品转化为人类知识很重要。

\subsection{开放问题}
\label{subsec:evolution-open-problems}

第一个开放问题是评估器通用性。当前的成功集中在具有自动验证的领域：代码、数学、游戏、调度、内核和基准测试的智能体。许多重要的AI任务涉及模糊性、人类偏好、长期社会效应或难以评分的安全约束。将这些领域的自进化扩展需要更丰富的评估器、不确定性估计和人-在-环治理。

第二个开放问题是理论理解。LLM驱动的变异不是一个简单的随机过程。它以预训练、提示词设计、归档内容和模型解码为条件。经典演化算法的收敛性分析不直接适用。我们需要语义变异、归档诱导偏差和评估器噪声的模型。没有理论，系统设计仍然是经验性的和脆弱的。

第三个开放问题是数据污染和基准耗竭。随着系统反复优化公开基准，基准变得信息量减少。自进化可以通过发现基准的怪癖来消耗它。这对于软件工程任务尤为严重，因为代码库、测试和解决方案可能泄漏到模型训练数据中。动态基准、私有测试和具有可验证奖励的任务生成是部分补救措施。

第四个开放问题是自修改代码的安全性。演化程序可能消耗过多资源、访问非预期文件或学会操纵其评估器。运行时自我修改引发额外担忧，因为修改机制本身可以被更改。实用系统需要沙箱化、最小权限执行、来源追踪、策略检查和紧急回滚。安全性应作为一等指标进行评估，而非作为实现细节处理。

第五个开放问题是多目标对齐。实际部署关心准确率、成本、延迟、可维护性、隐私、公平性和鲁棒性。标量基准分数无法编码所有这些。MAP-Elites和Pareto归档提供了一条路径，但描述符选择很困难。如果描述符质量差，归档可能保留不相关的多样性而遗漏重要的多样性。

第六个开放问题是自进化与权重学习之间的关系。OPRO、FunSearch、AlphaEvolve、DGM和SE-Agent可以通过外部搜索改进冻结模型系统。Absolute Zero通过自我对弈RL更新模型策略。未来系统可能结合两者：外部归档发现制品，模型定期在成功的谱系上训练或蒸馏。挑战在于保留可验证的增益而不蒸馏评估器技巧或降低多样性。

第七个开放问题是科学信用分配。当演化程序解决数学问题时，哪个组件值得信赖：基础模型、提示词、评估器、归档、人类问题表述还是计算预算？这不仅是哲学问题。它影响可复现性、知识产权以及研究人员如何解释发现。感知谱系的归档和详细的实验日志对于可信声明是必要的。

\subsection{综合}
\label{subsec:evolution-synthesis}

演化式代码与算法发现标志着AI自进化从孤立生成到累积搜索的转变。OPRO表明候选和分数的自然语言历史可以将LLM变成黑盒优化器。FunSearch表明候选可以是可执行程序，其演化行为可以产生数学发现。AlphaEvolve表明同一思想可以通过模型集成、自动化评估器和质量-多样性归档扩展到更广泛的算法和基础设施领域。OpenEvolve和CodeEvolve表明社区正在将这些思想转化为可重用框架。SE-Agent表明演化操作可以作用于推理轨迹。ADAS、DGM、G{\"o}del Agent、Agent Symbolic Learning、Absolute Zero和SelfEvolve表明智能体自身可以成为演化的基底。

统一方程很简单：
\begin{equation}
  \text{自进化} = \text{生成} + \text{评估} + \text{记忆} + \text{选择} + \text{约束下的变异}.
\end{equation}
每一项都是必要的。没有评估的生成是想象。没有记忆的评估是反复试错。没有选择的记忆是无结构的日志。没有变异的选择是停滞。没有约束的变异是不安全的。本章综述的系统在如何实例化各项上有所不同，但它们都依赖相同的闭环原则。

对于未来的综述和基准工作，最重要的建议是报告整个演化生态。论文应指定制品表示、模型路由、提示构造、归档策略、多样性描述符、评估器细节、验证划分、计算预算、失败率和回滚策略。没有这些细节，自进化结果难以比较且容易夸大。有了它们，该领域可以从令人印象深刻的演示走向可靠、可审计和可重用的自我改进AI系统。


\subsection{逐系统分析笔记}
\label{subsec:evolution-system-notes}

以下笔记通过将共同分析视角应用于每个系统来扩展比较。它们被纳入是因为演化自改进方法通常用相似词汇描述，即使其工程假设差异很大。对于每个系统，关键问题是：什么在演化，什么反馈被信任，什么记忆被保留，什么失败模式占主导。

\paragraph{OPRO。} OPRO演化指令字符串和序列化的优化变量。其可信反馈信号是评分历史，这意味着外循环的可靠性仅与分配分数、错误或奖励的机制一样可靠。在本章使用的自进化分类中，OPRO不应被理解为独立的模型能力，而应被理解为生成器、评估器和归档之间的交互模式。生成器提供候选变异；评估器将行为转化为选择信号；归档决定哪些追踪存活到未来上下文中。这一视角很有用，因为它防止了一个常见误解：LLM不是在隔离中神奇地改进，而是在跨尝试积累证据的更大控制论系统中参与。

对于OPRO，核心工程问题是候选制品获得多少自由度。窄制品边界使验证更简单并降低系统改变自身测量过程的可能性。宽制品边界增加了发现定性新策略的机会，但也增加了沙箱化、来源和回归测试的负担。主导风险是提示过拟合和噪声评分器反馈。因此鲁棒实现应包括保留验证、制品谱系、失败摘要和回滚路径。在比较研究中，OPRO不仅应通过其最佳分数来评估，还应通过分数作为评估预算的函数、保留候选的多样性以及在重复种子下改进的稳定性来评估。

\paragraph{FunSearch。} FunSearch演化版本化的Python函数。其可信反馈信号是沙箱化的数学目标，这意味着外循环的可靠性仅与分配分数、错误或奖励的机制一样可靠。在本章使用的自进化分类中，FunSearch不应被理解为独立的模型能力，而应被理解为生成器、评估器和归档之间的交互模式。生成器提供候选变异；评估器将行为转化为选择信号；归档决定哪些追踪存活到未来上下文中。这一视角很有用，因为它防止了一个常见误解：LLM不是在隔离中神奇地改进，而是在跨尝试积累证据的更大控制论系统中参与。

对于FunSearch，核心工程问题是候选制品获得多少自由度。窄制品边界使验证更简单并降低系统改变自身测量过程的可能性。宽制品边界增加了发现定性新策略的机会，但也增加了沙箱化、来源和回归测试的负担。主导风险是评估器泄漏和过早收敛。因此鲁棒实现应包括保留验证、制品谱系、失败摘要和回滚路径。在比较研究中，FunSearch不仅应通过其最佳分数来评估，还应通过分数作为评估预算的函数、保留候选的多样性以及在重复种子下改进的稳定性来评估。

\paragraph{AlphaEvolve。} AlphaEvolve演化算法代码制品。其可信反馈信号是自动化程序和基础设施评估器，这意味着外循环的可靠性仅与分配分数、错误或奖励的机制一样可靠。在本章使用的自进化分类中，AlphaEvolve不应被理解为独立的模型能力，而应被理解为生成器、评估器和归档之间的交互模式。生成器提供候选变异；评估器将行为转化为选择信号；归档决定哪些追踪存活到未来上下文中。这一视角很有用，因为它防止了一个常见误解：LLM不是在隔离中神奇地改进，而是在跨尝试积累证据的更大控制论系统中参与。

对于AlphaEvolve，核心工程问题是候选制品获得多少自由度。窄制品边界使验证更简单并降低系统改变自身测量过程的可能性。宽制品边界增加了发现定性新策略的机会，但也增加了沙箱化、来源和回归测试的负担。主导风险是大型计算预算和有限的公开可复现性。因此鲁棒实现应包括保留验证、制品谱系、失败摘要和回滚路径。在比较研究中，AlphaEvolve不仅应通过其最佳分数来评估，还应通过分数作为评估预算的函数、保留候选的多样性以及在重复种子下改进的稳定性来评估。

\paragraph{OpenEvolve。} OpenEvolve演化用户定义的代码和提示词。其可信反馈信号是自定义评估器回调和制品，这意味着外循环的可靠性仅与分配分数、错误或奖励的机制一样可靠。在本章使用的自进化分类中，OpenEvolve不应被理解为独立的模型能力，而应被理解为生成器、评估器和归档之间的交互模式。生成器提供候选变异；评估器将行为转化为选择信号；归档决定哪些追踪存活到未来上下文中。这一视角很有用，因为它防止了一个常见误解：LLM不是在隔离中神奇地改进，而是在跨尝试积累证据的更大控制论系统中参与。

对于OpenEvolve，核心工程问题是候选制品获得多少自由度。窄制品边界使验证更简单并降低系统改变自身测量过程的可能性。宽制品边界增加了发现定性新策略的机会，但也增加了沙箱化、来源和回归测试的负担。主导风险是质量取决于从业者的评估器设计。因此鲁棒实现应包括保留验证、制品谱系、失败摘要和回滚路径。在比较研究中，OpenEvolve不仅应通过其最佳分数来评估，还应通过分数作为评估预算的函数、保留候选的多样性以及在重复种子下改进的稳定性来评估。

\paragraph{CodeEvolve。} CodeEvolve演化标记的代码区域和提示种群。其可信反馈信号是具有岛屿编排的沙箱指标，这意味着外循环的可靠性仅与分配分数、错误或奖励的机制一样可靠。在本章使用的自进化分类中，CodeEvolve不应被理解为独立的模型能力，而应被理解为生成器、评估器和归档之间的交互模式。生成器提供候选变异；评估器将行为转化为选择信号；归档决定哪些追踪存活到未来上下文中。这一视角很有用，因为它防止了一个常见误解：LLM不是在隔离中神奇地改进，而是在跨尝试积累证据的更大控制论系统中参与。

对于CodeEvolve，核心工程问题是候选制品获得多少自由度。窄制品边界使验证更简单并降低系统改变自身测量过程的可能性。宽制品边界增加了发现定性新策略的机会，但也增加了沙箱化、来源和回归测试的负担。主导风险是分布式执行复杂性。因此鲁棒实现应包括保留验证、制品谱系、失败摘要和回滚路径。在比较研究中，CodeEvolve不仅应通过其最佳分数来评估，还应通过分数作为评估预算的函数、保留候选的多样性以及在重复种子下改进的稳定性来评估。

\paragraph{SE-Agent。} SE-Agent演化多步推理轨迹。其可信反馈信号是SWE-bench Verified问题解决，这意味着外循环的可靠性仅与分配分数、错误或奖励的机制一样可靠。在本章使用的自进化分类中，SE-Agent不应被理解为独立的模型能力，而应被理解为生成器、评估器和归档之间的交互模式。生成器提供候选变异；评估器将行为转化为选择信号；归档决定哪些追踪存活到未来上下文中。这一视角很有用，因为它防止了一个常见误解：LLM不是在隔离中神奇地改进，而是在跨尝试积累证据的更大控制论系统中参与。

对于SE-Agent，核心工程问题是候选制品获得多少自由度。窄制品边界使验证更简单并降低系统改变自身测量过程的可能性。宽制品边界增加了发现定性新策略的机会，但也增加了沙箱化、来源和回归测试的负担。主导风险是多轨迹探索的成本。因此鲁棒实现应包括保留验证、制品谱系、失败摘要和回滚路径。在比较研究中，SE-Agent不仅应通过其最佳分数来评估，还应通过分数作为评估预算的函数、保留候选的多样性以及在重复种子下改进的稳定性来评估。

\paragraph{ADAS。} ADAS演化以Python程序编写的智能体架构。其可信反馈信号是智能体基准分数，这意味着外循环的可靠性仅与分配分数、错误或奖励的机制一样可靠。在本章使用的自进化分类中，ADAS不应被理解为独立的模型能力，而应被理解为生成器、评估器和归档之间的交互模式。生成器提供候选变异；评估器将行为转化为选择信号；归档决定哪些追踪存活到未来上下文中。这一视角很有用，因为它防止了一个常见误解：LLM不是在隔离中神奇地改进，而是在跨尝试积累证据的更大控制论系统中参与。

对于ADAS，核心工程问题是候选制品获得多少自由度。窄制品边界使验证更简单并降低系统改变自身测量过程的可能性。宽制品边界增加了发现定性新策略的机会，但也增加了沙箱化、来源和回归测试的负担。主导风险是庞大的图灵完备搜索空间。因此鲁棒实现应包括保留验证、制品谱系、失败摘要和回滚路径。在比较研究中，ADAS不仅应通过其最佳分数来评估，还应通过分数作为评估预算的函数、保留候选的多样性以及在重复种子下改进的稳定性来评估。

\paragraph{DGM。} DGM演化自修改编码智能体的归档。其可信反馈信号是编码基准改进，这意味着外循环的可靠性仅与分配分数、错误或奖励的机制一样可靠。在本章使用的自进化分类中，DGM不应被理解为独立的模型能力，而应被理解为生成器、评估器和归档之间的交互模式。生成器提供候选变异；评估器将行为转化为选择信号；归档决定哪些追踪存活到未来上下文中。这一视角很有用，因为它防止了一个常见误解：LLM不是在隔离中神奇地改进，而是在跨尝试积累证据的更大控制论系统中参与。

对于DGM，核心工程问题是候选制品获得多少自由度。窄制品边界使验证更简单并降低系统改变自身测量过程的可能性。宽制品边界增加了发现定性新策略的机会，但也增加了沙箱化、来源和回归测试的负担。主导风险是安全归档管理和昂贵评估。因此鲁棒实现应包括保留验证、制品谱系、失败摘要和回滚路径。在比较研究中，DGM不仅应通过其最佳分数来评估，还应通过分数作为评估预算的函数、保留候选的多样性以及在重复种子下改进的稳定性来评估。

\paragraph{Godel Agent。} Godel Agent演化运行时方法和行为逻辑。其可信反馈信号是猴子补丁后的验证，这意味着外循环的可靠性仅与分配分数、错误或奖励的机制一样可靠。在本章使用的自进化分类中，Godel Agent不应被理解为独立的模型能力，而应被理解为生成器、评估器和归档之间的交互模式。生成器提供候选变异；评估器将行为转化为选择信号；归档决定哪些追踪存活到未来上下文中。这一视角很有用，因为它防止了一个常见误解：LLM不是在隔离中神奇地改进，而是在跨尝试积累证据的更大控制论系统中参与。

对于Godel Agent，核心工程问题是候选制品获得多少自由度。窄制品边界使验证更简单并降低系统改变自身测量过程的可能性。宽制品边界增加了发现定性新策略的机会，但也增加了沙箱化、来源和回归测试的负担。主导风险是运行时不确定性和回滚安全。因此鲁棒实现应包括保留验证、制品谱系、失败摘要和回滚路径。在比较研究中，Godel Agent不仅应通过其最佳分数来评估，还应通过分数作为评估预算的函数、保留候选的多样性以及在重复种子下改进的稳定性来评估。

\paragraph{Agent Symbolic Learning。} Agent Symbolic Learning演化提示词、工具和流水线边。其可信反馈信号是语言损失和梯度，这意味着外循环的可靠性仅与分配分数、错误或奖励的机制一样可靠。在本章使用的自进化分类中，Agent Symbolic Learning不应被理解为独立的模型能力，而应被理解为生成器、评估器和归档之间的交互模式。生成器提供候选变异；评估器将行为转化为选择信号；归档决定哪些追踪存活到未来上下文中。这一视角很有用，因为它防止了一个常见误解：LLM不是在隔离中神奇地改进，而是在跨尝试积累证据的更大控制论系统中参与。

对于Agent Symbolic Learning，核心工程问题是候选制品获得多少自由度。窄制品边界使验证更简单并降低系统改变自身测量过程的可能性。宽制品边界增加了发现定性新策略的机会，但也增加了沙箱化、来源和回归测试的负担。主导风险是依赖强LLM生成的梯度。因此鲁棒实现应包括保留验证、制品谱系、失败摘要和回滚路径。在比较研究中，Agent Symbolic Learning不仅应通过其最佳分数来评估，还应通过分数作为评估预算的函数、保留候选的多样性以及在重复种子下改进的稳定性来评估。

\paragraph{Absolute Zero。} Absolute Zero演化自生成的任务和求解器策略。其可信反馈信号是代码执行器奖励，这意味着外循环的可靠性仅与分配分数、错误或奖励的机制一样可靠。在本章使用的自进化分类中，Absolute Zero不应被理解为独立的模型能力，而应被理解为生成器、评估器和归档之间的交互模式。生成器提供候选变异；评估器将行为转化为选择信号；归档决定哪些追踪存活到未来上下文中。这一视角很有用，因为它防止了一个常见误解：LLM不是在隔离中神奇地改进，而是在跨尝试积累证据的更大控制论系统中参与。

对于Absolute Zero，核心工程问题是候选制品获得多少自由度。窄制品边界使验证更简单并降低系统改变自身测量过程的可能性。宽制品边界增加了发现定性新策略的机会，但也增加了沙箱化、来源和回归测试的负担。主导风险是课程崩溃和验证器范围。因此鲁棒实现应包括保留验证、制品谱系、失败摘要和回滚路径。在比较研究中，Absolute Zero不仅应通过其最佳分数来评估，还应通过分数作为评估预算的函数、保留候选的多样性以及在重复种子下改进的稳定性来评估。

\paragraph{SelfEvolve。} SelfEvolve演化具有自生成知识的生成代码。其可信反馈信号是沙箱错误和断言，这意味着外循环的可靠性仅与分配分数、错误或奖励的机制一样可靠。在本章使用的自进化分类中，SelfEvolve不应被理解为独立的模型能力，而应被理解为生成器、评估器和归档之间的交互模式。生成器提供候选变异；评估器将行为转化为选择信号；归档决定哪些追踪存活到未来上下文中。这一视角很有用，因为它防止了一个常见误解：LLM不是在隔离中神奇地改进，而是在跨尝试积累证据的更大控制论系统中参与。

对于SelfEvolve，核心工程问题是候选制品获得多少自由度。窄制品边界使验证更简单并降低系统改变自身测量过程的可能性。宽制品边界增加了发现定性新策略的机会，但也增加了沙箱化、来源和回归测试的负担。主导风险是有限的纠正范围和手写提示词。因此鲁棒实现应包括保留验证、制品谱系、失败摘要和回滚路径。在比较研究中，SelfEvolve不仅应通过其最佳分数来评估，还应通过分数作为评估预算的函数、保留候选的多样性以及在重复种子下改进的稳定性来评估。


\subsection{可复现演化发现的方法论检查清单}
\label{subsec:evolution-methodological-checklist}

一篇可复现的演化发现论文应该像机器学习论文描述数据、模型和优化器那样仔细描述搜索过程。第一个必需项是\emph{候选模式}。论文应声明候选是提示词、完整程序、补丁、标记函数、轨迹、工具图还是模型生成的任务。它还应声明周围系统的哪些部分是不可变的。没有这个边界，读者无法区分算法改进和基准线束的意外更改。对于代码制品，模式应包括允许的导入、允许的文件写入、时间限制、内存限制、依赖版本以及候选是否可以观察评估器内部。

第二个项是\emph{提议分布}。在LLM引导的演化中，提议分布不是简单的高斯变异。它是模型选择、解码温度、提示结构、父代采样、上下文压缩和输出解析的产物。因此可复现性报告应包括模型名称、模型版本或日期、采样参数、每个提示的样本数、重试策略、提示模板和任何后处理规则。如果系统使用双模型策略（如廉价模型用于广度、更强模型用于深度），应报告每个模型何时被调用以及模型路由决策如何做出。如果失败被摘要回模型，摘要模板应该是方法的一部分，因为它改变了未来的变异。

第三个项是\emph{归档策略}。自进化系统由它记住什么来定义。它是否存储每个候选、仅有效候选、仅精英或也包括失败？它是否保持谱系？候选是否按行为签名、MAP-Elites单元格、提示嵌入、任务领域或手工设计的描述符聚类？父代选择偏好近期候选、高适应度候选、新颖候选还是未被充分探索的生态位？这些选择可以主导性能。仅最佳归档近似爬山；岛屿归档近似分布式演化搜索；MAP-Elites归档近似景观照明。论文应报告归档大小、修剪规则、重置规则以及验证结果是否可以覆盖训练分数。

第四个项是\emph{评估器契约}。评估器应描述为具有输入、输出、失败模式和信任假设的接口。对于数学程序搜索，评估器可能返回确定性分数。对于代码智能体，它可能返回基准问题的通过/失败。对于智能体轨迹，它可能涉及执行环境和评判器。契约应定义什么算作无效、如何处理超时、是否可能有部分信用以及随机结果是否被平均。隐藏测试和公开测试应分开。如果评估器昂贵，论文应报告运行了多少次评估以及在评分前丢弃了多少。

第五个项是\emph{预算核算}。演化系统可能仅仅因为消耗更多尝试而看起来比单次基线强得多。公平比较不要求所有用例的预算相同，但确实需要透明的预算。作者应报告词元计数、模型调用计数、候选计数、有效候选计数、评估器调用、挂钟时间、硬件和可能时的货币成本。最有信息的图表通常是性能与评估器预算的关系，而非仅仅最终分数。一个快速达到中等分数的系统可能优于一个仅在大量搜索后才达到更高分数的系统。

第六个项是\emph{回归测量}。候选不应仅因在单一批次上改进一个标量而被接受。系统应测量它是否保留了先前能力、是否泛化到保留任务以及是否在对抗或随机测试下保持安全。在代码演化中，回归套件应包括功能正确性、性能、确定性和资源使用。在智能体演化中，回归套件应包括工具使用安全、适当时的拒绝行为、记忆一致性和对提示扰动的鲁棒性。接受门控应是显式的。对于多目标系统，论文应声明它使用词典排序、标量化、Pareto支配还是生态位特定精英。

第七个项是\emph{谱系报告}。没有谱系的最终制品难以信任。谱系报告应包括父代标识符、变异提示、评估结果和接受原因。对于科学发现，它应包括揭示思想如何涌现的中间制品。这使人类研究人员能够检查系统是发现了真正的机制还是利用了评估器怪癖。谱系对调试也很有用：如果后续世代退化，归档可以识别退化进入的分支。

第八个项是\emph{负面证据}。演化发现论文通常强调成功的分支，但失败的分支具有科学价值。它们揭示了评估器拒绝了什么、模型倾向于误解什么以及哪些搜索算子是无效的。简洁的失败分类可能比一长串原始失败程序更有用。典型类别包括语法错误、语义类型错误、超时失败、测试过拟合、数值不稳定性、不确定性、依赖误用、不安全文件访问和奖励黑客。报告失败分布帮助其他研究人员复现搜索并设计更好的过滤器。

第九个项是\emph{人工干预核算}。一些系统声称自主自进化，但人类可能在运行期间调整提示词、修复基础设施、选择基准或手动检查候选。这些干预不一定是坏事；它们对于早期研究通常是必要的。但它们应该被报告。自主搜索、人工策划搜索和事后人工解释之间的清晰区分防止了过度声明。对于部署，可接受的人工干预水平将取决于风险：科学探索可能容忍人工策划，而自主生产代码修改需要更严格的自动化和监督。

第十个项是\emph{发布制品}。可复现的发布应包括最终制品、种子制品、提示模板、评估器代码、基准划分、配置文件和归档日志。如果专有模型阻止精确复现，发布仍然可以提供足够结构以使用其他模型进行近似复现。FunSearch、OpenEvolve、CodeEvolve和SE-Agent等开源系统之所以有价值，是因为它们暴露了这些周围基础设施，而不仅仅是因为它们发布了巧妙的提示词。

\subsection{演化自改进循环的概念分类}
\label{subsec:evolution-taxonomy-loops}

本章中的系统可以按演化介入的深度来组织。最浅的级别是\emph{指令演化}。制品是自然语言提示词，评估器测量任务性能。OPRO是代表性案例。指令演化易于部署，因为它不需要代码执行，但受限于提示词的表达力和LLM行为在提示词变化下的稳定性。它最适合固定模型已经具有潜在能力并需要更好引出的任务。

下一个级别是\emph{程序片段演化}。制品是函数或标记的代码区域。FunSearch是最干净的示例。此级别具有强可解释性，因为演化片段可以被阅读和分析。它还具有强遏制性，因为评估器可以限制片段的输入和输出。限制是周围的脚手架必须已经很好地定义问题。如果关键创新需要更改脚手架，单函数搜索可能太窄。

第三个级别是\emph{代码库演化}。制品可能包括多个函数、文件或算法模块。AlphaEvolve、OpenEvolve和CodeEvolve更接近这一级别。代码库演化可以发现更大的设计更改并优化真实系统，但它需要构建系统、集成测试、性能线束、依赖管理和更强的沙箱化。它还需要更好的信用分配，因为性能更改可能来自多次编辑之间的交互。

第四个级别是\emph{轨迹演化}。制品是思想、动作、工具调用或中间计划的序列。SE-Agent是代表性的。当最终代码更改不如智能体到达它的过程重要时，轨迹演化很有吸引力。修订或重组的轨迹可以教会智能体如何处理类似问题。风险是轨迹质量可能在很大程度上取决于基础模型和环境状态，使复现更加困难。

第五个级别是\emph{架构演化}。制品是智能体设计、工作流或多智能体组织。ADAS通过让元智能体编写智能体程序来搜索这一级别。DGM在归档中演化智能体实现。架构演化可以发现新的推理和工具使用模式，但它具有最广的搜索空间和最难的评估问题。好的架构可能跨任务鲁棒，而坏的架构可能以微妙的方式过拟合到基准怪癖。

第六个级别是\emph{自指运行时演化}。制品是智能体自身在执行时的操作逻辑。G{\"o}del Agent说明了这一级别。运行时演化很强大，因为适应可以在线发生，但它需要强回滚和遏制。系统必须确保变更机制不会破坏评估机制。这是安全担忧最严重的级别。

第七个级别是\emph{通过自生成数据的策略演化}。Absolute Zero代表这一级别。被改进的制品不再仅仅是外部代码或提示词；它是模型策略本身，通过自我对弈任务和可验证奖励训练。此级别可以产生更持久的改进，因为知识可能进入权重，但它也提高了训练的成本和风险。它需要奖励设计、课程控制和防止自生成分布崩溃的保障措施。

这些级别不是互斥的。未来的自进化系统可能使用OPRO风格的提示词优化来改进其变异提示词、FunSearch风格的程序演化来发现启发式方法、MAP-Elites来保持代码多样性、SE-Agent风格的轨迹重组来改进问题解决追踪、ADAS风格的架构搜索来重新设计智能体，以及Absolute-Zero风格的自我对弈来训练底层模型。研究挑战是如何组合这些循环而不创建不受控的反馈周期。

\subsection{对更广泛综述的启示}
\label{subsec:evolution-implications-survey}

对于AI自进化的更广泛主题，演化式代码与算法发现提供了最具体的证据，表明AI系统可以在常规监督训练之外积累改进。证据是具体的，因为制品是可检查的且评估器通常是可执行的。提示词可以被阅读，函数可以被运行，补丁可以被测试，归档可以被审计。这使得该领域成为可能否则变得推测性的讨论的有用锚点。

本章还表明自进化应作为系统学科来研究。LLM很重要，但它不是整个系统。评估器、归档、沙箱、提示编译器、模型路由器、基准线束和日志基础设施同样重要。在许多情况下，创新在于这些组件如何连接。FunSearch的版本化提示格式、AlphaEvolve的质量-多样性归档、DGM的开放式智能体集合和Agent Symbolic Learning的语言梯度图都是系统级贡献。

第二个启示是可验证性将塑造近期自进化的版图。代码、数学、类定理任务、游戏、调度、芯片设计和内核优化可能会进展更快，因为它们提供可靠的反馈。谈判、教育、医疗和社会规划等领域将需要更仔细的评估器设计和治理。该领域应抵制在无条件限定的情况下从可验证领域泛化到所有智能形式的诱惑。

第三个启示是多样性既是安全机制也是能力机制。它是能力机制，因为多样化归档保留踏脚石和替代策略。它是安全机制，因为多样性使得比较制品和检测向评估器技巧的可疑收敛变得更容易。如果每个分支独立发现相同的捷径，研究人员应该问捷径是否反映了预期目标。如果不同分支通过不同机制解决任务，人类获得更多洞见。

第四个启示是自进化改变了人类研究人员的角色。人类可能花更少时间编写最终启发式方法，更多时间设计搜索空间、评估器、约束和审计过程。这不是科学责任的减少。这是从直接构造到生态系统设计的转变。由此产生的发现质量取决于生态系统是否奖励真正的进步。

第五个启示是发表规范需要演进。自进化结果不能仅由最终分数和方法图来完整描述。它需要归档摘要、预算报告、评估器规格、谱系证据和失败分析。没有这些，自主改进的声明难以解释。有了它们，社区可以在可靠性、效率和通用性上比较系统，而不仅仅在最大基准分数上。

总之，演化式代码与算法发现是当前AI自进化研究的运营核心。它展示了LLM如何在嵌入具有记忆和验证的循环时充当优化器、变异算子、重组器、调试器和元设计师。该领域的下一步是使这些循环更可复现、更安全、更预算感知和更通用，同时保持使程序演化具有科学价值的可解释性。
